\chapter{Legislação --- Segurança da Informação e Proteção de Dados}
\label{cap:legislacao}

\begin{edital}
Lei 12.527/2011 (LAI) caps.\ I--V + Decretos 7.724 e 7.845; Lei 12.737/2012
(Delitos Informáticos) art.\ 2º; Lei 12.965/2014 (Marco Civil) cap.\ II Seção
I e cap.\ III Seções I e II; Lei 13.709/2018 (LGPD) caps.\ I, II, III, IV,
VII, VIII, IX.
\end{edital}
\peso{5 questões, peso 1. A FGV adora prazo, competência e a distinção
controlador $\times$ operador. Ela cita o número e a data da lei no
enunciado --- não é decoreba de número, é entender competências, prazos e
papéis.}

\begin{pegadinha}[Atualização crítica --- Marco Civil, art.\ 19, pós-STF jun/2025]
Em \textbf{26/06/2025} o STF mudou a regra de responsabilidade dos provedores
por conteúdo de terceiros --- é a atualização legislativa mais provável de
cair em 2026 e tem seção própria mais adiante neste capítulo (ver
Seção \ref{sec:marco-civil}). Não decore a redação antiga do art.\ 19 como se
ainda vigesse: hoje a regra geral é \textbf{notice-and-takedown}.
\end{pegadinha}

\section{LGPD (Lei 13.709/2018)}

\begin{itemize}
\item \textbf{Dado pessoal} (art.\ 5º, I) é informação relacionada a pessoa
natural \textbf{identificada ou identificável} --- nome, endereço, telefone,
e-mail, placa, CPF. \textbf{Dado pessoal sensível} (art.\ 5º, II) é um rol
\textbf{fechado e menor}: origem racial ou étnica, convicção religiosa,
opinião política, filiação a sindicato ou a organização de caráter religioso,
filosófico ou político, dado referente à \textbf{saúde} ou à \textbf{vida
sexual}, dado \textbf{genético} ou \textbf{biométrico} vinculado a uma pessoa
natural. Sensível tem regime de tratamento \textbf{mais restrito} (art.\ 11),
não proibição.
\item \textbf{Controlador} decide sobre o tratamento (finalidade e meios);
\textbf{operador} trata em nome do controlador. \textbf{Encarregado (DPO)} é
o canal com titulares e ANPD.
\item \textbf{Bases legais (art.\ 7º):} consentimento \textbf{não} é a
única; há \textbf{10} bases (cumprimento de obrigação legal, execução de
contrato, legítimo interesse, proteção da vida, tutela da saúde, políticas
públicas, etc.).
\item \textbf{Direitos do titular (art.\ 18):} confirmação, acesso,
correção, anonimização/bloqueio/eliminação, portabilidade, informação sobre
compartilhamento, revogação do consentimento.
\item \textbf{ANPD} (autarquia) \textbf{fiscaliza e sanciona}; o
\textbf{CNPD} (Conselho) é \textbf{consultivo} (propõe diretrizes, sugere
ações à ANPD). Não confunda os papéis.
\item \textbf{Sanções:} advertência; \textbf{multa simples de até 2\% do
faturamento} no Brasil no último exercício, excluídos tributos, limitada a
\textbf{R\$ 50 milhões por infração}; publicização; bloqueio/eliminação dos
dados. Exige processo administrativo.
\item \textbf{LGPD aplicada a IA:} explicabilidade em decisão automatizada
(art.\ 20), anonimização, \emph{scraping} --- os três já cobrados no TJ-RJ.
\end{itemize}

\subsection{Princípios (art.\ 6º) --- o trio que a FGV confunde}

São \textbf{dez} princípios: finalidade, adequação, necessidade, livre
acesso, qualidade dos dados, transparência, segurança, prevenção, não
discriminação, e responsabilização e prestação de contas. Três deles são o
alvo preferido da banca, porque parecem sinônimos e não são:

\begin{center}
\begin{tabular}{@{}p{2.4cm}p{9cm}@{}}
\toprule
\textbf{Princípio} & \textbf{O que exige} \\
\midrule
\textbf{Finalidade} & propósitos \textbf{legítimos, específicos e explícitos},
\textbf{informados ao titular}, sem tratamento posterior incompatível \\
\textbf{Adequação} & \textbf{compatibilidade} do tratamento com as finalidades
informadas, \textbf{conforme o contexto} \\
\textbf{Necessidade} & tratamento \textbf{limitado ao mínimo} indispensável
para atingir a finalidade (pertinente, proporcional, não excessivo) \\
\bottomrule
\end{tabular}
\end{center}

A âncora é a palavra-chave: \textbf{finalidade} = \emph{para quê} (e avisou);
\textbf{adequação} = \emph{combina com} o que foi avisado; \textbf{necessidade}
= \emph{o mínimo}.

\begin{pegadinha}
No TJ-RJ a FGV descreveu literalmente ``o tratamento seja \textbf{compatível
com os fins informados ao titular, de acordo com o contexto}'' e ofereceu
\emph{finalidade}, \emph{prevenção}, \emph{adequação}, \emph{necessidade} e
\emph{transparência} --- gabarito \textbf{adequação}. Quem tinha decorado só
``finalidade'' errou, porque a palavra ``fins'' está no enunciado de
propósito. Leia até o fim: quem fala em \emph{compatibilidade} e
\emph{contexto} é a adequação; quem fala em \emph{informar o propósito} é a
finalidade; quem fala em \emph{mínimo} é a necessidade.
\end{pegadinha}

\begin{pegadinha}
Troca \textbf{controlador $\times$ operador}: controlador \textbf{decide}
sobre o tratamento; operador \textbf{executa} em nome do controlador. Troca
\textbf{ANPD $\times$ CNPD}: a ANPD é autarquia federal de natureza especial
que fiscaliza e aplica sanção; o CNPD é órgão consultivo --- apenas
propõe/sugere e dissemina (no Dataprev 2024 a alternativa certa era o CNPD
``sugerir ações à ANPD''). Dizer que consentimento é a única base legal da
LGPD é falso --- há 10 bases.
\end{pegadinha}

\section{Marco Civil (Lei 12.965/2014)}
\label{sec:marco-civil}

\subsection{Guarda de registros}

\begin{itemize}
\item \textbf{Registros de conexão:} guarda obrigatória por \textbf{1 ano}
(provedor de conexão).
\item \textbf{Registros de acesso a aplicações:} \textbf{6 meses} (provedor
de aplicação).
\item Princípios: neutralidade de rede, privacidade, liberdade de
expressão.
\end{itemize}

\textbf{Neutralidade de rede (art.\ 9º).} O responsável pela transmissão,
comutação ou roteamento tem o dever de \textbf{tratar de forma isonômica
quaisquer pacotes de dados, sem distinção por conteúdo, origem e destino,
serviço, terminal ou aplicação}. Ou seja: o provedor não pode degradar, nem
priorizar, nem bloquear tráfego por causa \emph{do que} ele carrega ou
\emph{de quem} ele vem. A discriminação só cabe nas \textbf{exceções do
§1º}: requisitos técnicos indispensáveis à prestação adequada dos serviços e
priorização de \textbf{serviços de emergência} --- e ainda assim mediante
regulamentação.

\begin{pegadinha}
Os distratores de neutralidade descrevem justamente o que a lei proíbe, com
cara de eficiência operacional: ``priorizar o tráfego de quem paga mais'',
``bloquear aplicações de concorrentes'', ``monitorar o conteúdo dos pacotes
para fins comerciais'', ``reduzir a velocidade de vídeo no horário de pico
para aliviar a rede''. Esse último é o mais perigoso: soa técnico, mas é
discriminação \textbf{por aplicação}, que não está entre as exceções do §1º.
\end{pegadinha}

\subsection{Art.\ 19 --- responsabilidade dos provedores (atualização crítica pós-STF, jun/2025)}

O art.\ 19 do Marco Civil trata de quando uma plataforma/provedor responde
civilmente por conteúdo \textbf{gerado por terceiros} (posts, comentários,
vídeos de usuários). Até 2025, a redação original condicionava essa
responsabilidade ao descumprimento de \textbf{ordem judicial específica} de
remoção. Isso \textbf{mudou}.

\begin{conceito}[O que o STF decidiu --- Temas 987 e 533, 26/06/2025]
Em \textbf{26/06/2025}, o STF julgou dois recursos extraordinários com
repercussão geral (\textbf{Tema 987} e \textbf{Tema 533}) e, por
\textbf{8 votos a 3}, declarou o \textbf{art.\ 19 parcial e
progressivamente inconstitucional}. A tese fixada vale \textbf{a partir do
julgamento} e continua valendo \textbf{até que o Congresso legisle} sobre o
tema especificamente --- ou seja, é a regra vigente \textbf{agora}, não uma
proposta.
\end{conceito}

\begin{center}
\begin{tabular}{@{}p{3cm}p{4.6cm}p{4.6cm}@{}}
\toprule
& \textbf{Antes (redação original)} & \textbf{Agora (pós-STF, 26/06/2025)} \\
\midrule
Regra geral & provedor só responde se descumprir \textbf{ordem judicial específica} de remoção & regra geral virou \textbf{notice-and-takedown}: \textbf{notificação extrajudicial} já pode responsabilizar o provedor \\
Conteúdo íntimo não consentido (art.\ 21) & já bastava notificação (exceção histórica) & continua bastando notificação (a exceção virou a regra geral) \\
Crimes contra a honra (calúnia, injúria, difamação) & exigia ordem judicial & \textbf{continua exigindo ordem judicial} --- é a única categoria que preserva a regra antiga \\
Conteúdo de risco grave (crime contra o Estado/ordem democrática, terrorismo, incitação ao racismo, conteúdo que ponha em risco a vida) & tratamento igual ao geral & exige atuação \textbf{proativa e imediata} do provedor, sob dever de cuidado --- não espera nem notificação \\
\bottomrule
\end{tabular}
\end{center}

Em resumo: o Marco Civil deixou de exigir ordem judicial como regra --- a
\textbf{notificação extrajudicial} (do usuário prejudicado, de terceiro ou de
autoridade) já é suficiente para obrigar o provedor a agir na maioria dos
casos. A \textbf{ordem judicial permanece obrigatória apenas para os crimes
contra a honra}, e para conteúdo de risco grave o provedor deve agir de
forma \textbf{proativa}, sem esperar sequer notificação.

\begin{pegadinha}[Como a FGV vai usar a lei antiga como distrator]
A partir de 2026, espere que a banca ofereça a redação \textbf{original} do
art.\ 19 como se ainda fosse a regra --- e essa alternativa vai \textbf{soar
correta} para quem estudou por material desatualizado. Reconheça o padrão:

\begin{itemize}
\item \textbf{Distrator típico:} ``o provedor de aplicações de internet só
pode ser responsabilizado civilmente por danos decorrentes de conteúdo
gerado por terceiros se, após \textbf{ordem judicial específica}, não tomar
as providências para tornar indisponível o conteúdo apontado como
infringente'' --- isso era a regra \textbf{até 25/06/2025}. Hoje é
\textbf{a exceção} (vale só para crimes contra a honra), não mais a regra
geral. Se a alternativa não mencionar nenhuma ressalva de data ou de tipo de
crime, desconfie.
\item \textbf{Distrator inverso (superextrapolação):} dizer que ``o STF
extinguiu totalmente a exigência de ordem judicial'' --- também
\textbf{falso}: a ordem judicial \textbf{continua obrigatória} para calúnia,
injúria e difamação. A regra não é absoluta em nenhuma das duas direções.
\item \textbf{Distrator de terminologia:} chamar o novo regime de
``censura prévia'' --- é \textbf{notice-and-takedown} (reação a notificação
sobre conteúdo já publicado), não filtragem antes da publicação.
\item \textbf{Distrator de data/instrumento:} atribuir a mudança a uma
\textbf{lei nova do Congresso} --- foi uma \textbf{decisão do STF em controle
de constitucionalidade} (Temas 987 e 533), com efeito \textbf{até que} o
Congresso legisle sobre o tema; não confundir instância nem data.
\end{itemize}
\end{pegadinha}

\section{LAI (Lei 12.527/2011) --- prazos de sigilo}

\begin{center}
\begin{tabular}{@{}p{4cm}p{6cm}@{}}
\toprule
\textbf{Classificação} & \textbf{Prazo máximo de restrição} \\
\midrule
Reservada & 5 anos \\
Secreta & 15 anos \\
Ultrassecreta & 25 anos (prorrogável \textbf{uma vez}) \\
\bottomrule
\end{tabular}
\end{center}

Regra geral: \textbf{publicidade é a regra; sigilo é exceção.} Prazos contam
da \textbf{data de produção} da informação. Informação sobre
\textbf{violação de direitos humanos} por agente público \textbf{não} pode
ser objeto de restrição de acesso. \textbf{Transparência ativa} (o órgão
publica de ofício, ex.: portal) $\times$ \textbf{passiva} (o cidadão pede via
e-SIC).

\textbf{Desclassificação e reavaliação.} A classificação \textbf{não é
definitiva nem discricionária}:

\begin{itemize}
\item \textbf{Findo o prazo} (ou consumado o evento que define o termo
final), a informação torna-se \textbf{automaticamente} de acesso público ---
\textbf{não} é preciso procedimento próprio nem decisão específica para
liberar.
\item A lei permite \textbf{desclassificar} \textbf{e também reduzir o
prazo}. A reavaliação cabe à \textbf{autoridade classificadora ou a
autoridade hierarquicamente superior}, \textbf{de ofício ou mediante
provocação}.
\item A decisão de classificar é vinculada a balizas legais (grau, prazo
máximo, competência por autoridade), não é escolha livre do agente.
\end{itemize}

\textbf{Decretos que o edital cita} (regulamentam a LAI no Executivo
federal):

\begin{itemize}
\item \textbf{Decreto 7.724/2012:} regulamenta a LAI --- procedimentos de
acesso, e-SIC, prazos de resposta ao pedido (\textbf{20 dias}, prorrogáveis
por \textbf{10}), recursos, transparência ativa, competências de
classificação.
\item \textbf{Decreto 7.845/2012:} procedimentos de \textbf{credenciamento
de segurança} e tratamento de \textbf{informação classificada} (custódia,
acesso, controle).
\end{itemize}

\begin{pegadinha}
Inverter prazos (LAI, Marco Civil) e competências (ANPD $\times$ CNPD).

Na Dataprev 2024 os distratores da LAI foram exatamente estes: trocar os
nomes dos graus (``ultrassigilosas, sigilosas ou reservadas'' --- é
\textbf{ultrassecreta, secreta, reservada}) e dizer que a decisão de
classificar é \textbf{discricionária}, ``porquanto a lei não prevê balizas'';
afirmar que, findo o prazo, ``são necessários procedimento próprio e decisão
específica'' (é \textbf{automático}); e afirmar que a lei ``permite a
desclassificação, \textbf{mas não} a redução de prazo'' (permite as duas).
\end{pegadinha}

\section{Delitos Informáticos (Lei 12.737/2012, art.\ 154-A do CP)}

Crime de \textbf{invasão de dispositivo informático}. \textbf{Independe} de
o dispositivo estar conectado à rede (a conexão não é elementar do tipo).
Aumento de pena conforme o sujeito passivo (autoridades) e o prejuízo. (A
pena exata foi alterada pela Lei 14.155/2021 --- foque no conceito.)

\begin{pegadinha}
Dizer que a invasão do art.\ 154-A exige o aparelho conectado à rede é
\textbf{falso} --- o Dataprev derrubou exatamente essa alternativa.
\end{pegadinha}

\begin{jacaiu}
\textbf{Em prova real da FGV:} \textbf{ANPD $\times$ CNPD} --- o gabarito foi
justamente ``o CNPD tem atribuição de \emph{sugerir ações} a serem
realizadas pela ANPD''; \textbf{classificação na LAI} (nomes dos graus,
decisão \emph{vinculada} a balizas, acesso automático findo o prazo,
desclassificação \emph{e} redução de prazo); \textbf{sanções administrativas
da LGPD} (parâmetros de dosimetria, destino da arrecadação das multas);
\textbf{sanções do Marco Civil} (advertência com prazo para adoção de
medidas); \textbf{invasão de dispositivo} do art.\ 154-A --- não exige
conexão à rede, e o aumento de pena varia com o sujeito passivo ---
\textbf{Dataprev 2024}. \textbf{Princípio da adequação} (o enunciado descreve
``compatível com os fins informados, de acordo com o contexto'');
\textbf{explicabilidade da decisão automatizada}, com a redação do art.\ 20;
\textbf{IA generativa $\times$ LGPD}, incluindo raspagem da web ---
\textbf{TJ-RJ}.

\textbf{No nosso banco} (previsto pelo edital, ainda não visto na amostra de
provas): controlador $\times$ operador; as dez bases legais; dado pessoal
sensível; os prazos numéricos da LAI (5/15/25); transparência ativa $\times$
passiva; guarda de registros no Marco Civil (1 ano / 6 meses); neutralidade
de rede; e todo o regime do art.\ 19 pós-STF de 26/06/2025
(notice-and-takedown, exceção dos crimes contra a honra, dever proativo,
repercussão geral).

Rode \texttt{./quiz.py legislacao}.
\end{jacaiu}

\section{Alta probabilidade / pesquisa extra}

\begin{itemize}
\item \textbf{GDPR} (regulamento europeu) aparece como paralelo à LGPD.
\item Atenção a novas decisões do STF/ANPD em 2026 --- legislação é o bloco
que mais muda; confira antes da prova.
\end{itemize}

\begin{comosair}
LAI, prazos máximos contados da produção da informação: reservada 5 anos,
secreta 15, ultrassecreta 25. Só a ultrassecreta prorroga, uma única vez,
por igual período. LGPD, multa simples: até 2\% do faturamento no Brasil no
último exercício, excluídos tributos, limitada a R\$ 50 milhões por
infração. Marco Civil art.\ 19 (STF, 26/06/2025): regra geral virou
notice-and-takedown --- notificação extrajudicial já basta, exceto para
crimes contra a honra (ordem judicial). Guarde: ANPD fiscaliza/sanciona;
CNPD só aconselha.
\end{comosair}
